\chapter{Integration: Jarvis-PLOT}
\label{ch:jarvisplot}

\newthought{\JPlot\ turns scan output into figures.} It is a separate, optional tool, but \Jarvis\
is built to hand off to it smoothly: one command turns your task card into a plotting
configuration you can render. This chapter covers that hand-off; \JPlot\ has its own
documentation for the full plotting language.

\section{What Jarvis-PLOT is}
\label{sec:plot-what}

With \JPlot\ you do not write plotting code---you \emph{describe a figure in \textsc{yaml}} and the
\cmd{jplot} command loads the data, runs any transforms, and renders the image. If you know
Matplotlib, the model is familiar: pick a \code{method} (\code{scatter}, \code{contourf},
\ldots) and pass styling, but declaratively. Install it separately (Chapter~\ref{ch:installation}):

\begin{shellbox}
python3 -m pip install -U JarvisPLOT
\end{shellbox}

\section{The hand-off from a scan}
\label{sec:plot-handoff}

\Jarvis\ does not draw posterior figures itself; it \emph{generates a \JPlot\ configuration} from
your task card. Run the scan, then ask \Jarvis\ for a plot config with \cli{-{}-plot}
(Chapter~\ref{ch:cli}), and render it with \cmd{jplot}:

\begin{shellbox}
Jarvis bin/my_scan.yaml --plot
jplot images/my_scan/my_scan.yaml
\end{shellbox}

\noindent \cli{-{}-plot} writes a starter \textsc{yaml} into \file{images/<scan>/}; you then edit
that file to choose figures, columns, and styling, and re-run \cmd{jplot} to redraw. Because the
config points at the run's \textsc{csv}/\textsc{hdf5} (Chapter~\ref{ch:outputs}), you can also
snapshot a still-running scan with \cli{-{}-convert} and plot the partial results.

\begin{jtip}
\cmd{jplot -{}-parse-data} inspects the run's \textsc{csv}/\textsc{hdf5}, sanitises column names,
and writes the discovered columns back into the plotting \textsc{yaml}---a quick way to see exactly
which observables are available to plot.
\end{jtip}

\section{The workflow flowchart}
\label{sec:plot-flowchart}

\JPlot\ also renders the workflow flowchart. \Jarvis\ always writes the semantic
\file{images/<scan>/flowchart.json}; if \JPlot\ is installed, \Jarvis\ asks it to render
\file{flowchart.png} (Chapter~\ref{ch:workflow-model}). \cli{-{}-skip-draw-flowchart} skips the
picture while still exporting the \textsc{json}.

\section{Where the boundary is}
\label{sec:plot-boundary}

\Jarvis\ owns the data and the semantic config; \JPlot\ owns layout, rendering, and styling. So
this manual stops at the hand-off: for the plotting \textsc{yaml} schema, figure types, transforms,
and the gallery, see the \JPlot\ documentation.

\section{Summary}
\label{sec:plot-summary}

Run a scan, generate a config with \cli{-{}-plot}, and render with \cmd{jplot}; \JPlot\ also draws
the workflow \file{flowchart.png}. \Jarvis\ produces the data and the starting config; \JPlot\
produces the figures.
